\section{Theories and their classifying topoi}\label{sec:theories}

So far the decks were presheaf topoi given by a shape. In Caramello's programme, the most interesting decks come from logic: a theory and its classifying topos (\S\ref{sec:sites}). For theories of presheaf type, the arches of \S\ref{sec:arches} still apply, with $C$ the opposite of the category of finitely presented models.

\subsection{Algebraic theories satisfy the weak excluded middle}

\begin{theorem}\label{th:horn}
Let $\mathbb T$ be a Horn (cartesian) theory, for instance the theory of groups, of abelian groups, of rings, of modules over a fixed ring, of lattices, or of sets with no structure. Then:
\begin{enumerate}
\item its classifying topos $\Set[\mathbb T]$ is De Morgan;
\item $\Set[\mathbb T]$ is Boolean iff every homomorphism between finitely presented models of $\mathbb T$ is an isomorphism.
\end{enumerate}
\end{theorem}

\inwords{} The internal logic of the \q{universe of all groups} (or rings, or modules) satisfies $\neg p\vee\neg\neg p$, though not $p\vee\neg p$.

\gen{} Deck: $\Set[\mathbb T]\simeq[\fp\mathbb T\text{-mod},\Set] = \PSh(C)$ with $C = (\fp\mathbb T\text{-mod})^{\op}$ (Example~\ref{ex:classifying}). Arch B: a cospan in $C$ is a span $B\leftarrow A\rightarrow B'$ of finitely presented models, and a completing square is a pair of homomorphisms $B\to D\leftarrow B'$ agreeing on $A$: the \emph{amalgamation property}. For a Horn theory it holds, because finitely presented models have pushouts (glue the presentations: generators of $B$ and $B'$, relations of both, and the identification of the images of the generators of $A$). Arch A: $C$ is a groupoid iff all homomorphisms of finitely presented models are invertible.

\direct{} A proof without topoi would require the description of subobjects in $\Set[\mathbb T]$, which is beyond this text; here the bridge is the proof. For the algebra: given $A\to B$ and $A\to B'$, with presentations $B = \langle X\mid R\rangle$, $B' = \langle X'\mid R'\rangle$, and $A$ generated by $Y$, the model $D = \langle X\sqcup X'\mid R\cup R'\cup\{\alpha(y) = \alpha'(y) : y\in Y\}\rangle$ is finitely presented and receives compatible maps from $B$ and $B'$.\qed

\begin{example}
\begin{enumerate}
\item The theory with no axioms, on one sort: $\fp$ models are finite sets, $C = \mathsf{Fin}^{\op}$, which satisfies Ore ($\mathsf{Fin}$ has pushouts). Not Boolean: the map $\emptyset\to\{\ast\}$ is not invertible.
\item The theory of an \emph{inhabited} set ($\top\vdash\exists x\,\top$): its finitely presented models are the non-empty finite sets, and its classifying topos is $[\mathsf{Fin}_{>0},\Set]$. Pushouts of non-empty sets are non-empty: De Morgan.
\item The theory of a \emph{decidable} set (a set with a relation $\neq$ complementary to $=$) has as models finite sets with injections: $[\mathsf{FI},\Set]$, De Morgan by Theorem~\ref{th:combi}(3).
\end{enumerate}
The last two are not Horn theories but are of presheaf type, and the same arch applies.
\end{example}

\status{} Known, in a stronger form: every regular theory is De Morgan \cite[Thm~3.5]{caramello-demorgan}, and a theory of presheaf type is De Morgan iff its finitely presented models have the amalgamation property \cite[Prop.~14]{caramello-chamoun}.

\subsection{Torsors}

\begin{theorem}\label{th:torsors}
For groups $G$ and $H$, the theory of $G$-torsors and the theory of $H$-torsors are Morita equivalent (have equivalent classifying topoi) iff $G\cong H$.
\end{theorem}

\inwords{} Two groups have \q{the same torsors in every topos} only if they are isomorphic.

\gen{} Deck: the classifying topos of $G$-torsors is $\PSh(\B G)$. Arch C: $\PSh(\B G)\simeq\PSh(\B H)$ iff $\Kar(\B G)\simeq\Kar(\B H)$; a group has no idempotent other than $1$, so $\Kar(\B G)\simeq\B G$, and $\B G\simeq\B H$ iff $G\cong H$ (an equivalence between one-object categories is fully faithful on the single object, hence an isomorphism of the groups of arrows).

\leanbox{\lean{ClassifyingTopoi.isClassifyingSite_genericTorsor}, \lean{ClassifyingTopoi.moritaEquivalent_torsorTheory_iff}.}

\status{} Standard.

\subsection{Quotients and topologies}

Bridge~\ref{br:sub} is the arch that Caramello's talk develops most \cite[slides 38--54]{caramello-mines}. The quotients of a theory $\mathbb T$ (theories with the same language and more axioms) form a lattice, which is isomorphic to the lattice of Grothendieck topologies on a classifying site, ordered by inclusion \cite[ch.~3]{caramello-tst}. Computing in the lattice of topologies, where there are explicit formulas for the topology generated by a family of sieves and for meets and implications \cite[slides 50--51]{caramello-mines}, gives theorems about theories that would be hard to find syntactically:
\begin{itemize}
\item a \emph{deduction theorem} for geometric logic: if $\psi$ is provable in $\mathbb T$ plus the axiom $\varphi$ (for sentences $\varphi$, $\psi$), then $\varphi\vdash\psi$ is provable in $\mathbb T$ \cite[slide 52]{caramello-mines};
\item an explicit axiomatisation of the \emph{meet} of the theory of local rings and the theory of integral domains, inside the lattice of quotients of the theory of commutative rings: the sequents $0 = 1\vdash\bot$, and, for polynomials $P_s, G_i, H_j$ in $\mathbb Z[x_1,\dots,x_n]$ such that the $P_s$ and $G_i$ generate the unit ideal and $\prod_j H_j\in(P_1,\dots,P_m)$, the sequents
\[
\textstyle\bigwedge_s P_s(\vec x) = 0\ \vdash_{\vec x}\ \bigvee_i\exists y\,(G_i(\vec x)\,y = 1)\ \vee\ \bigvee_j H_j(\vec x) = 0
\]
\cite[\S18]{caramello-lattices}.
\end{itemize}
These are beyond the level we aim at, and we do not reproduce the proofs. In the library, Bridge~\ref{br:sub} for classifying sites is \lean{ClassifyingTopoi.IsClassifyingSite.exists_quotient_theory} and its converse \lean{ClassifyingTopoi.IsClassifyingSite.exists_quotient_topology}; the matching of quotients of Morita-equivalent theories is \lean{ClassifyingTopoi.MoritaEquivalent.exists_orderIso_quotients}.

\begin{exercise}\label{exo:inhabited}
The theory of an inhabited set is the quotient of the theory of a set by the single axiom $\top\vdash\exists x\,\top$. Its classifying topos $[\mathsf{Fin}_{>0},\Set]$ is therefore a subtopos of $[\mathsf{Fin},\Set]$. By Bridge~\ref{br:sub} it corresponds to a Grothendieck topology on $\mathsf{Fin}^{\op}$. Which sieve on the object $\emptyset$ must be covering? (Hint: the axiom says that the model presented by \q{no generators} maps to the model presented by \q{one generator}.)
\end{exercise}
